\begin{frame}{기저율 오류를 1000명 세기로}
  \small
  \begin{tabular}{@{}lcc@{}}
    \toprule
     & 병 걸림 (1명) & 건강 (999명) \\
    \midrule
    검사 \textbf{양성} & $1\times0.99=$ \textbf{0.99} & $999\times0.01=$ \textbf{9.99} \\
    검사 음성 & 0.01 & 989.01 \\
    \bottomrule
  \end{tabular}
  \vskip0.8em
  \begin{itemize}
    \item ``양성''으로 좁히면 총 10.98명, 그중 병든 0.99명 ---
      $0.99/10.98 \approx 0.090$, 확률 계산과 \textbf{완전히 같은 답}.
    \item ``검사 정확도가 99\%인데 왜 사후확률이 9\%인가'' ---
      \textbf{가짜 양성과 진짜 양성의 비율}로 설명.
    \item 스팸 예의 확률 트리와 \textbf{동일한 구조}로 검증.
  \end{itemize}
\end{frame}
